% Definimos el estilo del documento
\documentclass[12pt,a4paper,spanish]{book}
% Utilizamos el paquete para utilizar espa~nol
\usepackage{babel}
% Utilizamos un paquete para gestionar los acentos
% y las e~nes
\usepackage[latin1]{inputenc}
% Utilizamos el paquete para gestionar imagenes jpg
\usepackage{graphicx}
% Definimos la zona de la pagina ocupada por el texto
\oddsidemargin -1.0cm
\headsep -2.4cm
\textwidth=18.5cm
\textheight=26cm
%Empieza el documento
\begin{document}
% Definimos titulo, autor, fecha.
\title{TITULO DE MI PROYECTO}
\author{Mar\'ia P\'erez}
\date{1 de octubre del 2015}
% Generamos titulo e indice de contenidos
\maketitle
\tableofcontents
% Definimos una primera pagina para los agradecimientos
\newpage
\thispagestyle{empty}
\section*{Agradecimientos}
Aqu\'{i} ponemos los agradecimientos
% Empezamos capitulos
\chapter{Introducci\'on}
2
Aqui empieza la introducci\'on
\chapter{Estado del arte}
Aqui empieza el capitulo sobre estado del arte\\
EJEMPLOS DE ECUACIONES\\
Las expresiones matem\'{a}ticas son
sencillas de escribir en el propio texto:
$a=\sum_{i=1}^{i=\infty} x_i^{n+1}$
y deben ser escritas entre d\'{o}lares. \\
Los super\'{\i}ndices
se obtienen con \^{},
$x^3 y^{\alpha + \beta}$, mientras que los
sub\'{\i}ndices son con \_  pudiendo
combinarlos.\\
Para la f\'{o}rmula  centrada
$$z^{2+\alpha}_{n+k}.$$
\begin{equation}
\label{integral}
f(x)=\int{f'(x)dx}+C
\end{equation}

Para centrar f\'omulas matem\'aticas tambi\'en podemos escribirlas entre corchetes. 
\[
x=\frac{a_2 x^2 + a_1 x + a_0}{1+2z^3}, \quad
x+y^{2n+2}=\sqrt{b^2-4ac}
\]
En las referencias \cite{La86} y \cite{Ro93} se encuentra una
descripci\'on en profundidad de las caracter\'{\i}sticas de
LaTex. La figura 1 que est\'a en la p\'agina 2, la tabla 1, el teorema
de Pitagoras y la f\'ormula 2 son ejemplos
de como utilizar las etiquetas.


\begin{thebibliography}{1}
\bibitem{La86} Leslie Lamport.
{\em LaTex : A document Preparation System}.
Addison-Wesley, 1986.
\bibitem{Ro93} Christian Rolland.
{\em LaTex guide pratique}.
Addison-Wesley, 1993.
\end{thebibliography}
% Termina el documento
\end{document}